\documentclass[10pt,a4paper]{amsart}
\usepackage[margin=19mm]{geometry}
\usepackage{amsmath,amssymb,mathtools}
\usepackage[scr=rsfs,cal=euler]{mathalfa}
\usepackage{xfrac}
\usepackage[unicode,hidelinks,bookmarks=false]{hyperref}
\newcommand{\ie}{i.e.\ }
\newcommand{\cf}{cf.\ }
\newcommand{\resp}{respectively\ }
\newcommand{\Subsection}{\subsection}
\providecommand{\nobreakdash}{\nobreak\hbox{-}}
\setlength{\emergencystretch}{2em}
\multlinegap0pt
\usepackage{bm}
\usepackage{bbm}
\usepackage{dsfont}
\usepackage{xspace}
\usepackage{cancel}
\usepackage{mathtools}
\usepackage{amsthm}
\makeatletter
\@ifundefined{theorem}{\newtheorem{theorem}{Theorem}}{}
\@ifundefined{lemma}{\newtheorem{lemma}[theorem]{Lemma}}{}
\makeatother
\newcommand{\eps}{\varepsilon}
\title[Local Second-Order Bounds for Aggregation-Order Variation in]{Local Second-Order Bounds for Aggregation-Order Variation in Density Fusion}
\author{Ratan Bahadur Thapa}
\date{Source: arXiv:2608.01420v2 (CC BY 4.0)}
\begin{document}
\maketitle

\noindent\emph{Source and licence.} Local Second-Order Bounds for Aggregation-Order Variation in Density Fusion. Version arXiv:2608.01420v2 (CC BY 4.0), distributed under the Creative Commons Attribution 4.0 International licence, as recorded in the arXiv licence metadata for this exact version and shown on the abstract page https://arxiv.org/abs/2608.01420v2. Full rights notice: licenses/arxiv-2608.01420-CC-BY-4.0.txt.\par\smallskip

\noindent\textit{Editorial scope.} Counted unit: the Lemma whose source label is \texttt{lem:root} in the article (source0.tex lines 356--358, source label \texttt{lem:root}), delivered together with the article's own complete original proof of it (source0.tex lines 360--392, one \texttt{proof} environment of 33 lines). No mathematics is written, repaired or re-derived: statement and proof are the article's own text line for line. Required context retained uncounted: lem:chart (lines 345--347). Every definition, display and label of those lines is retained unchanged. Omitted from this package and not redistributed: 1--344, 348--355, 359--359, 393--1791. Nothing inside a retained excerpt is omitted or altered: every retained body is delivered exactly as the article's lines are written, so no omission receipt of any kind accompanies this package. Citations: the retained excerpts carry no citation command at all, so no bibliography entry of the article is transcribed and no reference is redistributed. Reference adaptation: none was needed and none was made; every reference inside the delivered unit resolves inside it, so no unresolved reference or citation is produced. Printed slips kept literal: no printed word, symbol, hypothesis or number of the counted statement or of its proof was altered. Layout and class: the article's own class is \texttt{article} with packages including graphicx, multirow, amsmath, amssymb, amsfonts, amsthm, mathrsfs, xcolor, textcomp, manyfoot, booktabs, listings, xspace, geometry; the wrapper is the lane's pinned \texttt{amsart} editorial wrapper at 10pt on A4, which restates the theorem-like environments the retained text needs and adds the article's own macro definitions that the retained text needs (\texttt{\textbackslash eps}), with extra wrapper packages (bm, bbm, dsfont, xspace, cancel, mathtools). The delivered numbering is the unit's own. Assets: no retained span contains a figure environment, a graphics-inclusion command, a PDF-page inclusion, a table or a TikZ picture; no image file is copied into this package, the package declares no asset dependency and no image file is redistributed with it. Licence: version arXiv:2608.01420v2 of this article is distributed under the Creative Commons Attribution 4.0 International licence, as recorded in the arXiv licence metadata for this exact version and shown on the abstract page https://arxiv.org/abs/2608.01420v2. A full rights notice is delivered at licenses/arxiv-2608.01420-CC-BY-4.0.txt. Characters: every retained excerpt is pure ASCII, so the delivered engine, XeLaTeX with its default Latin Modern text font, reports no missing character.
\begin{lemma}\label{lem:chart}
For sufficiently small $|\eps|$, all segment densities $(1-t)r^\eps+ts^\eps$, $0\le t\le1$, are positive and uniformly comparable with $p$. In addition, $r^\eps/((1-t)r^\eps+ts^\eps)$ and $s^\eps/((1-t)r^\eps+ts^\eps)$ converge uniformly to $1$.
\end{lemma}

\begin{lemma}\label{lem:root}
There exist $\eta>0$ and $\eps_0>0$ such that, for $0<|\eps|<\eps_0$, the balancing equation has exactly one solution in $(t_0-\eta,t_0+\eta)$. The solution converges to $t_0$ as $\eps\to0$.
\end{lemma}

\begin{proof}[Proof.]
Let
\[
\Phi_\eps(t)=aD_f(r^\eps\|(1-t)r^\eps+ts^\eps)-bD_f(s^\eps\|(1-t)r^\eps+ts^\eps).
\]
The local Taylor expansion of the two divergences gives the uniform expansion, for $t$ in a neighborhood of $t_0$,
\[
\eps^{-2}\Phi_\eps(t)= \Phi_2(t)+\eps\Phi_3(t)+o(\eps),
\]
where $\Phi_3$ is bounded uniformly on that neighborhood and
\[
\Phi_2(t)=A_2J\{at^2-b(1-t)^2\}.
\]
The corresponding derivative expansion is
$\eps^{-2}\partial_t\Phi_\eps(t)= \Phi_2'(t)+o(1)$
uniformly near $t_0$. The bounded-ratio assumptions and the $C^3$ expansion of $f$ near $1$ justify the uniform expansion of both $\Phi_\eps$ and $\partial_t\Phi_\eps$ on that neighborhood. Indeed, if $p_t^\eps=(1-t)r^\eps+ts^\eps$, then differentiating the exact integrand gives
\[
\partial_t\left[p_t^\eps f\left(\frac{r^\eps}{p_t^\eps}\right)\right]
=(s^\eps-r^\eps)\left\{f\left(\frac{r^\eps}{p_t^\eps}\right)
-\frac{r^\eps}{p_t^\eps}f'\left(\frac{r^\eps}{p_t^\eps}\right)\right\},
\]
and the same identity with $s^\eps$ in place of $r^\eps$. Lemma~\ref{lem:chart}, bounded density ratios, and the $C^3$ expansion of $f$ near $1$ give the stated derivative expansion uniformly in $t$.
More explicitly, as $z\to0$,
\[
f'(1+z)=f''(1)z+\frac12f'''(1)z^2+o(z^2).
\]
Lemma~\ref{lem:chart} makes the endpoint-to-segment ratio increments converge uniformly to zero on the chosen compact neighborhood of $t_0$. A uniform Taylor bound and domination by an integrable multiple of $p$ then give the stated uniform derivative remainder.
The root $t_0$ is simple because
\[
\Phi_2'(t_0)=2A_2J\{at_0+b(1-t_0)\}>0.
\]
Choose $\eta>0$ with $\eta<\min\{t_0,1-t_0\}$ so that $\Phi_2'(t)>c>0$ for $|t-t_0|<\eta$. The stated derivative expansion gives $\partial_t\Phi_\eps(t)>0$ on the same interval for sufficiently small nonzero $\eps$. Therefore, $\Phi_\eps$ is strictly increasing on that interval. Since $\Phi_2(t_0-\eta)<0<\Phi_2(t_0+\eta)$ after shrinking $\eta$ if needed, the uniform expansion gives opposite signs for $\Phi_\eps$ at the two endpoints. The intermediate value theorem gives existence, and strict monotonicity gives uniqueness. Applying the endpoint sign argument on every smaller neighborhood of $t_0$ proves that $t_\eps\to t_0$.
\end{proof}

\end{document}
